% ============================================================
% Chapter 4: Methodology
% File: chapters/metodology.tex
% ============================================================

\section{Data}
\label{sec:data}

The dataset consists of annual GDP growth rates (expressed as percentage points)
for 20 Latin American countries over the period 1960--2024, sourced from the World
Bank Open Data repository. The countries included are: Argentina, Bolivia, Brazil,
Chile, Colombia, Costa Rica, Cuba, Dominican Republic, Ecuador, El Salvador,
Guatemala, Haiti, Honduras, Mexico, Nicaragua, Panama, Paraguay, Peru, Uruguay,
and Venezuela. This extends the dataset of \cite{padilla2020} by six additional
years (2019--2024), incorporating both the COVID-19 contraction of 2020 and the
heterogeneous post-pandemic recovery trajectories across the region.

Missing values are present for the year 1960 across most countries, and
additionally for El Salvador (1960--1965) and Cuba (1960--1971). All missing
values are handled prior to basis expansion via linear interpolation between
adjacent observed values, followed by forward-fill and backward-fill to address
leading and trailing gaps. This imputation is applied exclusively for the purpose
of constructing the functional representation; all reconstruction quality metrics
are evaluated only on observed, non-imputed data points.

\section{Basis Expansion Methods}
\label{sec:bases}

Functional Data Analysis requires that discrete observations
$\{(t_j,\, y_i(t_j))\}_{j=1}^{T}$ for each country $i$ be converted into a
continuous functional object $X_i(t)$ defined over
$\mathcal{T} = [1960,\, 2024]$. This is achieved by expressing each function
as a linear combination of basis elements:
\begin{equation}
    X_i(t) \;\approx\; \sum_{k=1}^{K} c_{ik}\,\phi_k(t),
    \label{eq:basis_expansion}
\end{equation}
where $\{\phi_k\}_{k=1}^{K}$ is the chosen basis system and $c_{ik}$ are the
estimated coefficients. The choice of basis is the central object of study of
this thesis. Three families are compared.

\subsection{B-spline Basis (Baseline)}
\label{subsec:bspline}

B-splines are piecewise polynomial functions of degree $p$ joined at interior
points called knots, with global continuity up to order $p-1$. They are the
standard basis in FDA for non-periodic data \cite{ramsay2005} and serve as
the baseline for comparison, replicating the approach of \cite{padilla2020}.

GDP growth trajectories are assumed to belong to the Sobolev space
$W_2^m(\mathcal{T})$, which requires that derivatives up to order $m$ be
square-integrable. The smoothing spline is obtained by minimizing:
\begin{equation}
    \sum_{j=1}^{T}\bigl(y_i(t_j) - X_i(t_j)\bigr)^2
    \;+\;
    \lambda\int_{\mathcal{T}}\bigl[X_i^{(m)}(t)\bigr]^2\,dt,
    \label{eq:smoothing_spline}
\end{equation}
where the smoothing parameter $\lambda$ controls the bias-variance tradeoff.
In this work $\lambda$ is chosen adaptively per country via a robust noise
estimate $\hat{\sigma}$, computed from the vector of first differences $\Delta y$
as:
\begin{equation}
    \hat{\sigma} \;=\;
    \frac{\mathrm{median}\!\left(|\Delta y - \mathrm{median}(\Delta y)|\right)}
         {0.6745\,\sqrt{2}},
    \label{eq:sigma_est}
\end{equation}
and the smoothing parameter is set to $\lambda = n\,\hat{\sigma}^2$, where $n$
is the number of observed points. This data-driven choice ensures that the
smoothing level is proportional to each country's intrinsic volatility rather
than fixed globally.

\subsection{Fourier Basis}
\label{subsec:fourier}

The Fourier basis consists of sinusoidal functions of increasing frequency and
is optimal for functions in $L^2(\mathcal{T})$ that exhibit periodic or
quasi-periodic behaviour. A Fourier approximation retaining $K$ frequency
components is:
\begin{equation}
    X_i(t) \;=\; a_0
    \;+\;\sum_{k=1}^{K}
    \Bigl[a_k\cos\!\Bigl(\tfrac{2\pi k\,t}{T}\Bigr)
    +\; b_k\sin\!\Bigl(\tfrac{2\pi k\,t}{T}\Bigr)\Bigr],
    \label{eq:fourier}
\end{equation}
implemented via the Discrete Fourier Transform (DFT) by retaining the $K$
lowest-frequency components and zeroing the remainder. Three levels are
compared: $K \in \{5,\,10,\,15\}$.

A known limitation of the Fourier basis for economic data is the
\emph{Gibbs phenomenon}: near discontinuities or structural breaks---such as
the 1982 Latin American Debt Crisis or the 2020 COVID-19 contraction---a
finite Fourier truncation produces spurious oscillations in the neighbourhood
of the jump. This is examined empirically in Chapter~\ref{chap:results}.

\subsection{Wavelet Basis}
\label{subsec:wavelets}

Wavelets provide an unconditional basis for Besov spaces
$B_{p,q}^{s}(\mathcal{T})$, which measure regularity in a way that tolerates
spatial inhomogeneity---that is, functions that are smooth over some time
intervals yet exhibit sharp transitions in others \cite{donoho1998}. A function
with a sparse wavelet representation belongs to a Besov ball of high regularity,
whereas a function requiring many wavelet coefficients to represent accurately
does not admit efficient approximation in either Sobolev or Fourier bases.

A multi-level discrete wavelet transform (DWT) decomposes a signal $y$ into
approximation and detail coefficients:
\begin{equation}
    y \;=\;
    \sum_{k} a_{J,k}\,\phi_{J,k}(t)
    \;+\;
    \sum_{j=1}^{J}\sum_{k} d_{j,k}\,\psi_{j,k}(t),
    \label{eq:dwt}
\end{equation}
where $\phi_{J,k}$ and $\psi_{j,k}$ are translates of the scaling function and
the mother wavelet at decomposition level $j$ and translation $k$, $J=4$ is the
number of decomposition levels used throughout, $a_{J,k}$ are the approximation
coefficients at the coarsest scale, and $d_{j,k}$ are the detail coefficients at
scale $j$.

The DWT is implemented as an explicit orthogonal matrix $\mathbf{W}$ constructed
with circular (periodic) boundary conditions, so that
$\mathbf{W}\mathbf{W}^{\!\top} = \mathbf{I}$ exactly. The inverse transform is
then $\mathbf{W}^{\!\top}$, guaranteeing machine-precision round-trip fidelity
independent of signal length. Three wavelet families of increasing smoothness
order are compared:

\begin{itemize}
    \item \textbf{Haar}: the simplest orthogonal wavelet, piecewise constant,
          with one vanishing moment. Suitable for detecting level shifts.
    \item \textbf{Daubechies~2 (DB2)}: the shortest smooth orthogonal wavelet,
          filter length 4, with one vanishing moment. Balances time and
          frequency localisation.
    \item \textbf{Symlet~4 (Sym4)}: a nearly symmetric variant of Daubechies
          wavelets, filter length 8, with three vanishing moments. Provides
          smoother reconstructions while maintaining compact support.
\end{itemize}

Denoising is achieved via \emph{soft thresholding} of the detail coefficients
following \cite{donoho1998}:
\begin{equation}
    \hat{d}_{j,k}
    \;=\;
    \mathrm{sgn}(d_{j,k})\,\max\!\bigl(|d_{j,k}| - \lambda,\;0\bigr),
    \label{eq:soft_threshold}
\end{equation}
where $\lambda \geq 0$ is the threshold parameter. The approximation
coefficients $a_{J,k}$ are never thresholded. Setting $\lambda = 0$ yields
lossless reconstruction; increasing $\lambda$ progressively removes fine-scale
detail, interpolating continuously between the raw-data representation and a
smooth approximation.

\section{Comparison Metrics}
\label{sec:metrics}

Three metrics are used to compare basis performance, each addressing a distinct
aspect of approximation quality.

\subsection{Reconstruction Error}
\label{subsec:rmse}

The root mean squared error (RMSE) between the observed values and the functional
reconstruction, evaluated exclusively on non-imputed observed time points:
\begin{equation}
    \mathrm{RMSE}_i
    \;=\;
    \sqrt{\frac{1}{|\mathcal{O}_i|}
    \sum_{t\,\in\,\mathcal{O}_i}
    \bigl(y_i(t) - \hat{X}_i(t)\bigr)^2},
    \label{eq:rmse}
\end{equation}
where $\mathcal{O}_i$ denotes the set of observed (non-missing) time points
for country $i$.

Because wavelets admit lossless reconstruction at $\lambda = 0$
(giving $\mathrm{RMSE} = 0$), a direct RMSE comparison across bases is only
meaningful at equivalent smoothing levels. The fair comparison strategy adopted
here is: for each country, find the wavelet threshold $\lambda^*$ such that
$\mathrm{RMSE}_i^{\mathrm{wav}}(\lambda^*) \approx \mathrm{RMSE}_i^{\mathrm{bsp}}$,
and then count the number of nonzero coefficients required. This directly tests
the Besov space hypothesis: at equivalent reconstruction quality, a sparser basis
requires fewer nonzero coefficients.

\subsection{Sparsity}
\label{subsec:sparsity_metric}

The compression ratio under soft thresholding, defined as the fraction of wavelet
coefficients set to zero:
\begin{equation}
    \mathrm{Sparsity}(\lambda)
    \;=\;
    1 - \frac{\|\hat{\mathbf{d}}(\lambda)\|_0}{|\mathbf{d}|},
    \label{eq:sparsity}
\end{equation}
where $|\mathbf{d}|$ denotes the total number of wavelet coefficients across all
decomposition levels.

\subsection{FPCA Variance Explained}
\label{subsec:fpca_metric}

Functional Principal Component Analysis (FPCA) is applied to the matrix of
functional reconstructions $\mathbf{X} \in \mathbb{R}^{N \times T}$
(where $N = 20$ countries and $T = 65$ time points). The centred matrix
$\widetilde{\mathbf{X}} = \mathbf{X} - \mathbf{1}\bar{X}^{\!\top}$ is
decomposed via the Singular Value Decomposition (SVD):
\begin{equation}
    \widetilde{\mathbf{X}} \;=\; \mathbf{U}\,\mathbf{S}\,\mathbf{V}^{\!\top},
    \label{eq:svd}
\end{equation}
where the columns of $\mathbf{V}$ are the functional principal components
(eigenfunctions of the empirical covariance operator), the diagonal entries of
$\mathbf{S}^2$ are proportional to the eigenvalues, and the rows of
$\mathbf{U}\mathbf{S}$ are the PC scores per country.

The percentage of variance explained by the $k$-th FPC is:
\begin{equation}
    \mathrm{VE}_k
    \;=\;
    \frac{s_k^2}{\displaystyle\sum_{j} s_j^2}
    \;\times\; 100\%,
    \label{eq:var_explained}
\end{equation}
and the number of FPCs required to reach 80\% cumulative variance is recorded
as $n_{80}$. A smaller $n_{80}$ indicates a more parsimonious functional
representation of the cross-country variation structure.

\section{Computational Implementation}
\label{sec:implementation}

All analyses are implemented in Python~3.12. B-spline fitting uses
\texttt{scipy.interpolate.UnivariateSpline}; Fourier expansions use
\texttt{scipy.fft}; wavelet transforms are implemented from first principles
using explicit orthogonal matrices, without relying on external wavelet
libraries, to ensure full transparency and reproducibility of the reconstruction
pipeline. Visualisations are produced with \texttt{matplotlib} and
\texttt{seaborn}. Complete source code is provided in Appendix~A.
